\section{Points and Coordinates in the Plane}

Before we introduce algebraic techniques, we need a place in which mathematical
objects can live and be seen. The real number line is useful, but it allows us
to move in only one direction. Many questions become clearer as soon as we can
move independently in two directions and draw what is happening. For this
reason, our first working environment will be the Cartesian plane.

The plane is simple enough to draw on paper, but already rich enough to contain
points, curves, regions, distances, directions, equations, and transformations.
Three-dimensional space will later extend many of the same ideas, but two
dimensions allow us to see the essential structure without hiding it behind a
complicated picture.

\subsection{Two Number Lines Crossing}

Take two copies of the real number line. Place one horizontally and the other
vertically so that they meet at their zero points.

\begin{center}
\begin{tikzpicture}[scale=0.9]
    \draw[axis] (-4.5,0) -- (4.7,0) node[right] {$x$};
    \draw[axis] (0,-3.5) -- (0,3.7) node[above] {$y$};

    \foreach \x in {-4,-3,-2,-1,1,2,3,4} {
        \draw (\x,0.08) -- (\x,-0.08) node[below=3pt] {\small $\x$};
    }
    \foreach \y in {-3,-2,-1,1,2,3} {
        \draw (0.08,\y) -- (-0.08,\y) node[left=3pt] {\small $\y$};
    }

    \node[below left=2pt] at (0,0) {$0$};
\end{tikzpicture}
\end{center}

The horizontal line is called the \emph{$x$-axis}, and the vertical line is
called the \emph{$y$-axis}. Their intersection is called the \emph{origin}. We
denote the origin by
\[
O=(0,0).
\]

The axes give us two independent measurements. The first tells us how far we
move horizontally from the origin. The second tells us how far we move
vertically.

\begin{definitionbox}
A point in the Cartesian plane is described by an ordered pair of real numbers
\[
(x,y).
\]
The first number is the \emph{$x$-coordinate}; the second is the
\emph{$y$-coordinate}.
\end{definitionbox}

The word \emph{ordered} is essential. The point $(2,5)$ is generally not the
same point as $(5,2)$. The first number always refers to horizontal position,
and the second always refers to vertical position.

\subsection{From Coordinates to a Point}

Consider
\[
P=(3,2).
\]
To locate this point, begin at the origin. Move three units to the right and two
units upward.

\begin{center}
\begin{tikzpicture}[scale=0.9]
    \draw[axis] (-1.2,0) -- (5.0,0) node[right] {$x$};
    \draw[axis] (0,-1.2) -- (0,4.0) node[above] {$y$};

    \foreach \x in {-1,1,2,3,4} {
        \draw (\x,0.07) -- (\x,-0.07) node[below=3pt] {\small $\x$};
    }
    \foreach \y in {-1,1,2,3} {
        \draw (0.07,\y) -- (-0.07,\y) node[left=3pt] {\small $\y$};
    }

    \draw[helper] (3,0) -- (3,2);
    \draw[helper] (0,2) -- (3,2);
    \fill (3,2) circle (2.2pt);
    \node[above right] at (3,2) {$P=(3,2)$};
\end{tikzpicture}
\end{center}

The dashed lines are auxiliary. They show how the coordinates are read from the
axes; they are not part of the point.

The signs of the coordinates encode direction:
\begin{itemize}
    \item positive $x$ means right of the $y$-axis;
    \item negative $x$ means left of the $y$-axis;
    \item positive $y$ means above the $x$-axis;
    \item negative $y$ means below the $x$-axis.
\end{itemize}

Thus geometry and arithmetic describe the same object in two languages:
\[
\text{point in the plane}
\quad\longleftrightarrow\quad
\text{ordered pair of real numbers}.
\]

\subsection{The Cartesian Plane as a Set}

Because every point is represented by two real numbers, the entire plane is
\[
\mathbb{R}^2
=
\{(x,y):x\in\mathbb{R},\ y\in\mathbb{R}\}.
\]

We read this as follows:

\begin{quote}
$\mathbb{R}^2$ is the set of all ordered pairs $(x,y)$ such that both $x$ and
$y$ are real numbers.
\end{quote}

The superscript $2$ does not mean that real numbers are squared. It records that
two real coordinates are required to specify one point.

The coordinate axes are themselves subsets of the plane:
\[
X=\{(x,y)\in\mathbb{R}^2:y=0\},
\qquad
Y=\{(x,y)\in\mathbb{R}^2:x=0\}.
\]
Their only common point is the origin:
\[
X\cap Y=\{(0,0)\}.
\]

\subsection{Conditions and Sets of Points}

A formula involving $x$ and $y$ may appear only as a condition. For example,
\[
y=2
\]
is a relation between two numbers. To define the corresponding geometric
object, we write the complete set:
\[
L=\{(x,y)\in\mathbb{R}^2:y=2\}.
\]

This is read as:

\begin{quote}
$L$ is the set of all points $(x,y)$ in the Cartesian plane such that $y=2$.
\end{quote}

The colon means ``such that.'' A vertical bar may be used instead:
\[
\{(x,y)\in\mathbb{R}^2:y=2\}
=
\{(x,y)\in\mathbb{R}^2\mid y=2\}.
\]

Geometrically, $L$ is a horizontal line.

\begin{center}
\begin{tikzpicture}[scale=0.8]
    \draw[axis] (-4.2,0) -- (4.4,0) node[right] {$x$};
    \draw[axis] (0,-1.3) -- (0,3.8) node[above] {$y$};
    \draw[subset] (-4,2) -- (4,2);
    \node[subset label,right] at (4,2) {$L$};
    \draw (0.08,2) -- (-0.08,2) node[left=3pt] {$2$};
\end{tikzpicture}
\end{center}

Similarly,
\[
M=\{(x,y)\in\mathbb{R}^2:x=-1\}
\]
is a vertical line.

\begin{center}
\begin{tikzpicture}[scale=0.8]
    \draw[axis] (-3.5,0) -- (3.8,0) node[right] {$x$};
    \draw[axis] (0,-3.1) -- (0,3.4) node[above] {$y$};
    \draw[subset] (-1,-3) -- (-1,3);
    \node[subset label,above] at (-1,3) {$M$};
    \draw (-1,0.08) -- (-1,-0.08) node[below=3pt] {$-1$};
\end{tikzpicture}
\end{center}

The distinction is fundamental:
\[
\boxed{\text{condition on coordinates}}
\qquad\text{versus}\qquad
\boxed{\text{set of all points satisfying the condition}}.
\]

\subsection{A Curved Set of Points}

Consider the condition
\[
x^2+y^2=4.
\]
The corresponding geometric object is
\[
C=\{(x,y)\in\mathbb{R}^2:x^2+y^2=4\}.
\]
We read this as ``the set of all points $(x,y)$ in the plane such that
$x^2+y^2=4$.'' Geometrically, $C$ is the circle of radius $2$ centred at the
origin.

\begin{center}
\begin{tikzpicture}[scale=0.9]
    \draw[axis] (-3.2,0) -- (3.4,0) node[right] {$x$};
    \draw[axis] (0,-3.2) -- (0,3.4) node[above] {$y$};
    \draw[subset] (0,0) circle (2);
    \fill (0,0) circle (1.7pt) node[below left] {$O$};
    \node[subset label,above right] at (1.4,1.4) {$C$};
\end{tikzpicture}
\end{center}

The equation is the condition. The set-builder expression states the ambient
space, the form of its elements, and the condition they must satisfy. The
diagram shows the geometric meaning of that set.

This gives one of our main working principles:
\[
\text{condition on coordinates}
\quad\longrightarrow\quad
\text{subset of }\mathbb{R}^2
\quad\longrightarrow\quad
\text{geometric interpretation}.
\]

The Cartesian plane therefore gives us a common language in which numerical
information, algebraic conditions, sets, and geometric objects can be compared
directly.
